\usepackage{tcolorbox}
\tcbuselibrary{theorems}
\tcbuselibrary{skins,xparse}
\tcbuselibrary{breakable}

\tcbset{
defstyle/.style={enhanced,
                  breakable,
                   colback=white,
                    colbacktitle=green!20,
                    coltitle=black,
                  fonttitle=\bfseries\upshape,
                  separator sign={.},
                  colframe=white,
                 leftrule=12mm,
                 overlay={\node[anchor=north west,rectangle,minimum width=1cm,minimum height=1cm,pattern={bricks},pattern color=green] at (frame.west) {}; }},
esestyle/.style={enhanced,
                  breakable,
                   colback=white,
                    colbacktitle=orange!20,
                    coltitle=black,
                  fonttitle=\bfseries\upshape,
                  separator sign={.},
                  colframe=white,
                 leftrule=12mm,
                 overlay={\node[anchor=north west,rectangle,minimum width=1cm,minimum height=1cm,pattern={sixpointed stars},pattern color=orange] at (frame.west) {}; }},
attstyle/.style={enhanced,
                  breakable,
                   colback=white,
                    colbacktitle=red!20,
                    coltitle=black,
                  fonttitle=\bfseries\upshape,
                  separator sign={.},
                  colframe=white,
                 leftrule=12mm,
                 overlay={\node[anchor=north west,rectangle,minimum width=1cm,minimum height=1cm,pattern={checkerboard},pattern color=red] at (frame.west) {}; }},        
ossstyle/.style={enhanced,
                  breakable,
                   colback=white,
                    colbacktitle=olive!20,
                    coltitle=black,
                  fonttitle=\bfseries\upshape,
                  separator sign={.},
                  colframe=white,
                 leftrule=12mm,
                  overlay={\node[anchor=north west,rectangle,minimum width=1cm,minimum height=1cm,pattern={crosshatch dots},pattern color=olive] at (frame.west) {}; }},           
procstyle/.style={enhanced,
                  breakable,
                   colback=white,
                    colbacktitle=blue!20,
                    coltitle=black,
                  fonttitle=\bfseries\upshape,
                  separator sign={.},
                  colframe=white,
                 leftrule=12mm,
                  overlay={\node[anchor=north west,rectangle,minimum width=1cm,minimum height=1cm,pattern={fivepointed stars},pattern color=blue!75!black] at (frame.west) {}; }},           
princstyle/.style={enhanced,
                  breakable,
                   colback=white,
                    colbacktitle=brown!20,
                    coltitle=black,
                  fonttitle=\bfseries\upshape,
                  separator sign={.},
                  colframe=white,
                 leftrule=12mm,
                 overlay={\node[anchor=north west,rectangle,minimum width=1cm,minimum height=1cm,pattern={north west lines},pattern color=brown] at (frame.west) {}; }}, 
teostyle/.style={enhanced,
                  breakable,
                   colback=white,
                    colbacktitle=cyan!20,
                    coltitle=black,
                  fonttitle=\bfseries\upshape,
                  separator sign={.},
                  colframe=white,
                 leftrule=12mm,
                 overlay={\node[anchor=north west,rectangle,minimum width=1cm,minimum height=1cm,pattern={grid},pattern color=cyan] at (frame.west) {}; }},                  
}
\newtcolorbox{lattention}{breakable,enhanced,arc=0mm,colback=gray!5,colframe=gray,leftrule=12mm,%
overlay={\node[anchor=north west,outer sep=2pt] at (frame.north west) {\includegraphics[width=8mm]{0022}}; }}
\newtcbtheorem [number within=chapter] {definizionet}{Definizione}{defstyle}{}
\newtcbtheorem [number within=chapter]{esempiot}{Esempio}{esestyle}{}
\newtcbtheorem [number within=chapter]{attenzionet}{Attenzione}{attstyle}{}
\newtcbtheorem [number within=chapter] {osservazionet}{Osservazione}{ossstyle}{}
\newtcbtheorem [number within=chapter] {procedurat}{Procedura}{procstyle}{}

\newtcbtheorem [number within=chapter] {principiot}{Principio}{princstyle}{}
\newtcbtheorem [number within=chapter] {teoremat}{Teorema}{teostyle}{}

\NewTColorBox[auto counter,number within=section]{exercise}{+O{}}{%
	enhanced,colframe=black,colback=white,coltitle=black,
	fonttitle=\bfseries,
	%underlay={%\begin{tcbclipinterior}
	%	\shade[inner color=green!80!yellow,outer color=yellow!10!white]
	%		(interior.north west) circle (2cm);
			%\draw[help lines,step=5mm,yellow!80!black,shift={(interior.north west)}]
			%(interior.south west) grid (interior.north east);
	%	\end{tcbclipinterior}},
		title={Esercizio~\thetcbcounter:},
		label={esercizio@\thetcbcounter},
		attach title to upper=\quad,
		after upper={\par\hfill\textcolor{black}%
			{\itshape Soluzione a pagina~\pageref{solution@\thetcbcounter}}},
		lowerbox=ignored,
		savelowerto=solutions/exercise-\thetcbcounter.tex,
		record={\string\solution{\thetcbcounter}{solutions/exercise-\thetcbcounter.tex}},
		#1
	}
	\NewTotalTColorBox{\solution}{mm}{%
		enhanced,colframe=black,colback=white,coltitle=black,
		fonttitle=\bfseries,
%		underlay={\begin{tcbclipinterior}
%				\shade[inner color=red!50!yellow,outer color=yellow!10!white]
%				(interior.north west) circle (2cm);
%				\draw[help lines,step=5mm,yellow!80!black,shift={(interior.north west)}]
%				(interior.south west) grid (interior.north east);
%			\end{tcbclipinterior}},
			title={Soluzione delle Esercizio~\ref{esercizio@#1} oa pagina~\pageref{esercizio@#1}:},
			phantomlabel={solution@#1},
			attach title to upper=\par,
		}{\input{#2}}
		\tcbset{no solution/.style={no recording,after upper=}}